% Fokus Nullstellen – Einheit 4 (quadratische-funktionen)
\clearpage
\setcounter{aufgabe}{5}
\einheitenkopf[e4][4 Nullstellen]{Einheit 4 · Nullstellen und Schnittpunkte berechnen}
\zweigzeile{Hier lernst du, die Nullstellen einer Parabel zu berechnen · ab Kl. 10 (am Gymnasium ab Kl. 9) · P10 oft · baut auf: Nr. 1–5}

\begin{aufgabe}{Ich erkenne, welche Gleichung die Nullstellen liefert.}
Schreibe die Gleichung auf, die du für die Nullstellen lösen musst. Löse sie nicht.
\begin{gleichungsraster}[1]
\gl{f(x) = (x-5)^2 - 4} & \gl{f(x) = x^2 + 3x - 10} \\
\gl{f(x) = -(x+1)^2 + 9} & \gl{f(x) = x^2 - 36} \\
\gl{f(x) = 3x^2 - 12x + 9} & \\
\end{gleichungsraster}
\end{aufgabe}

\verfahren{Nullstellen aus der Scheitelpunktform}

\begin{aufgabe}{Ich kann die Nullstellen aus der Scheitelpunktform berechnen.}
Berechne die Nullstellen.
\rechenplatz{4}
\begin{gleichungsraster}[4]
\gl{f(x) = (x-1)^2 - 4} & \gl{f(x) = (x-4)^2 - 1} \\
\gl{f(x) = (x-2)^2 - 9} & \gl{f(x) = (x-5)^2 - 16} \\
\end{gleichungsraster}
\end{aufgabe}

\begin{aufgabe}{Ich kann die Nullstellen aus der Scheitelpunktform berechnen – weiter.}
Berechne die Nullstellen.
\begin{gleichungsraster}[4]
\gl{f(x) = (x+2)^2 - 1} & \gl{f(x) = (x+4)^2 - 25} \\
\gl{f(x) = -(x-1)^2 + 16} & \gl{f(x) = -(x+3)^2 + 4} \\
\end{gleichungsraster}
\end{aufgabe}

\begin{aufgabe}{Ich kann die Nullstellen einer gestreckten Parabel berechnen.}
Berechne die Nullstellen.
\begin{gleichungsraster}[4]
\gl{f(x) = 2(x-3)^2 - 8} & \gl{f(x) = -3(x+1)^2 + 27} \\
\gl{f(x) = 0{,}5(x+3)^2 - 8} & \\
\end{gleichungsraster}
\end{aufgabe}

\begin{aufgabe}{Ich kann Nullstellen aus der Scheitelpunktform berechnen, die keine ganzen Zahlen sind.}
Berechne die Nullstellen. Gib sie genau und auf zwei Stellen nach dem Komma gerundet an.
\begin{gleichungsraster}[4]
\gl{f(x) = (x-3)^2 - 6} & \gl{f(x) = (x+1)^2 - 7} \\
\anweisung{Die Funktion $f$ mit $f(x) = x^2 + 8x + 13$ hat die Scheitelpunktform unten. Berechne die Nullstellen von $f$ genau und gerundet. (P10 2017 OS)}
\gl{f(x) = (x+4)^2 - 3} & \\
\end{gleichungsraster}
\end{aufgabe}

\begin{aufgabe}{Ich finde den Fehler bei Nullstellen aus der Scheitelpunktform.}
Finde den Fehler, benenne ihn und korrigiere die Rechnung.
\begin{teile}
\teil Lea berechnet die Nullstellen von $f(x) = (x-3)^2 - 25$.
\rechnung{0 &= (x-3)^2 - 25 &&\mid +25 \\ 25 &= (x-3)^2 &&\mid \sqrt{\phantom{x}} \\ 5 &= x - 3 &&\mid +3 \\ 8 &= x && \text{Nullstelle bei } x = 8}
\teil Ben berechnet die Nullstellen von $f(x) = (x+2)^2 - 9$.
\rechnung{0 &= (x+2)^2 - 9 &&\mid +9 \\ 9 &= (x+2)^2 \\ x &= 2 + 3 = 5 &&\text{oder}\quad x = 2 - 3 = -1}
\end{teile}
\end{aufgabe}

\begin{aufgabe}{Ich kann am Scheitel begründen, wie viele Nullstellen eine Parabel hat.}
Gib ohne Rechnung an, wie viele Nullstellen die Parabel hat, und begründe mit dem Scheitel.
\begin{teile}
\teil $f(x) = (x-4)^2 + 3$
\teil $f(x) = (x+2)^2 - 5$
\teil $f(x) = (x-6)^2$
\teil $f(x) = -(x-1)^2 + 2$
\teil $f(x) = -(x+3)^2 - 4$
\end{teile}
\end{aufgabe}

\verfahren{Nullstellen mit der p-q-Formel}

\begin{aufgabe}{Ich kann die Nullstellen mit der p-q-Formel berechnen.}
Berechne die Nullstellen.
\rechenplatz{4}
\begin{gleichungsraster}[4]
\gl{f(x) = x^2 - 6x + 8} & \gl{f(x) = x^2 + 4x + 3} \\
\gl{f(x) = x^2 - 8x + 12} & \gl{f(x) = x^2 + 10x + 9} \\
\end{gleichungsraster}
\end{aufgabe}

\begin{aufgabe}{Ich kann die Nullstellen mit der p-q-Formel berechnen – weiter.}
Berechne die Nullstellen.
\begin{gleichungsraster}[4]
\gl{f(x) = x^2 - 2x - 8} & \gl{f(x) = x^2 + 6x - 7} \\
\gl{f(x) = x^2 + x - 6} & \gl{f(x) = x^2 - 3x - 10} \\
\end{gleichungsraster}
\end{aufgabe}

\begin{aufgabe}{Ich kann die Nullstelle berechnen, wenn die Parabel die x-Achse nur berührt.}
Berechne die Nullstellen.
\begin{gleichungsraster}[4]
\gl{f(x) = x^2 - 10x + 25} & \gl{f(x) = x^2 + 8x + 16} \\
\end{gleichungsraster}
\end{aufgabe}

\begin{aufgabe}{Ich kann mit der p-q-Formel zeigen, dass eine Parabel keine Nullstelle hat.}
Berechne die Nullstellen.
\begin{gleichungsraster}[4]
\gl{f(x) = x^2 + 2x + 5} & \gl{f(x) = x^2 - 4x + 7} \\
\end{gleichungsraster}
\end{aufgabe}

\begin{aufgabe}{Ich kann die Nullstellen berechnen, wenn vor $x^2$ eine Zahl steht.}
Berechne die Nullstellen.
\begin{gleichungsraster}[4]
\gl{f(x) = 2x^2 + 4x - 30} & \gl{f(x) = -x^2 - 4x + 12} \\
\gl{f(x) = 0{,}5x^2 + x - 4} & \\
\anweisung{Berechne die Nullstellen der Funktion $f$. (P10 2020 OS)}
\gl{f(x) = 3x^2 + 9x + 6} & \\
\end{gleichungsraster}
\end{aufgabe}

\begin{aufgabe}{Ich kann mit der p-q-Formel Nullstellen berechnen, die keine ganzen Zahlen sind.}
Berechne die Nullstellen. Gib sie genau und auf zwei Stellen nach dem Komma gerundet an.
\begin{gleichungsraster}[4]
\gl{f(x) = x^2 - 4x + 1} & \gl{f(x) = x^2 + 2x - 4} \\
\anweisung{Gegeben ist die Funktion $f$ unten. Berechne die Nullstellen von $f$ genau und gerundet. (P10 2025 OS)}
\gl{f(x) = x^2 - 10x + 14} & \\
\end{gleichungsraster}
\end{aufgabe}

\begin{aufgabe}{Ich finde den Fehler bei Nullstellen mit der p-q-Formel.}
Finde den Fehler, benenne ihn und korrigiere die Rechnung.
\begin{teile}
\teil Mia berechnet die Nullstellen von $f(x) = x^2 - 8x + 7$.
\rechnung{0 &= x^2 - 8x + 7 && p = -8,\ q = 7 \\ x_{1,2} &= -4 \pm \sqrt{16 - 7} = -4 \pm 3 \\ x_1 &= -1, \quad x_2 = -7}
\teil Can berechnet die Nullstellen von $f(x) = 2x^2 + 12x + 10$.
\rechnung{0 &= 2x^2 + 12x + 10 && p = 12,\ q = 10 \\ x_{1,2} &= -6 \pm \sqrt{36 - 10} = -6 \pm \sqrt{26} \\ x_1 &\approx -0{,}90, \quad x_2 \approx -11{,}10}
\teil Jona berechnet die Nullstellen von $f(x) = x^2 + 4x - 21$.
\rechnung{0 &= x^2 + 4x - 21 && p = 4,\ q = -21 \\ x_{1,2} &= -2 \pm \sqrt{4 - 21} = -2 \pm \sqrt{-17} \\ & && \text{keine Nullstelle}}
\end{teile}
\end{aufgabe}
